\documentclass[a4paper,12pt,twoside]{memoir}

% Castellano
\usepackage[spanish,es-tabla,es-ucroman]{babel}
\selectlanguage{spanish}
\usepackage[T1]{fontenc}

% Fuente
\usepackage{lmodern}

% Imagenes
\usepackage{graphicx}
\graphicspath{{./img}}

% Utils
\usepackage{microtype}
\usepackage{placeins}
\usepackage[table]{xcolor}
\usepackage{tcolorbox}
\usepackage{subfig}
\usepackage[shortlabels]{enumitem}

% Enlaces
\usepackage[colorlinks]{hyperref}
\hypersetup{
    allcolors = {red}
}

% Capítulo
\chapterstyle{bianchi}

% Permitir mayúsculas con tilde en headers
\setlength{\headheight}{15.4pt}

% Variables
\title{Identificación de Parkinson por visión artificial}
\author{Catalin Andrei Cacuci}
\date{\today}
\newcommand{\thedni}{X7451927L}
\newcommand{\thetutor}{Alicia Olivares Gil}

%%%%%%%%%%%%%%%%%%%%%%%%%%%%%%%%%%%%%%%%%%%%%%%%%%%%%%%%%%%%%%%%%%%%%%%%%%%%%%%%
%                                   COMANDOS
%%%%%%%%%%%%%%%%%%%%%%%%%%%%%%%%%%%%%%%%%%%%%%%%%%%%%%%%%%%%%%%%%%%%%%%%%%%%%%%%

% Imagen de cabecera
\newcommand{\cabecera}{\noindent\includegraphics[width=\textwidth]{cabecera}}

% Página en blanco que no incrementa contador
\newcommand{\blankpage}{
    \null
    \thispagestyle{empty}
    \addtocounter{page}{-1}
    \newpage
}

% Comando para insertar una imagen en un lugar concreto.
% Los parámetros son:
% 1 --> Ruta absoluta/relativa de la figura
% 2 --> Texto a pie de figura
% 3 --> Tamaño en tanto por uno relativo al ancho de página
\newcommand{\imagen}[3]{
    \begin{figure}[!h]
        \centering
        \includegraphics[width=#3\textwidth]{#1}
        \caption{#2}\label{fig:#1}
    \end{figure}
    \FloatBarrier
}

% Comando para insertar una imagen sin posición.
% Los parámetros son:
% 1 --> Ruta absoluta/relativa de la figura
% 2 --> Texto a pie de figura
% 3 --> Tamaño en tanto por uno relativo al ancho de página
\newcommand{\imagenflotante}[3]{
    \begin{figure}
        \centering
        \includegraphics[width=#3\textwidth]{#1}
        \caption{#2}\label{fig:#1}
    \end{figure}
}

% El comando \figuraConPosicion nos permite insertar figuras comodamente, y
% utilizando siempre el mismo formato. Los parametros son:
% 1 --> Porcentaje del ancho de página que ocupará la figura (de 0 a 1)
% 2 --> Fichero de la imagen
% 3 --> Texto a pie de imagen
% 4 --> Etiqueta (label) para referencias
% 5 --> Opciones que queramos pasarle al \includegraphics
% 6 --> Opciones de posicionamiento a pasarle a \begin{figure}
\newcommand{\figuraConPosicion}[6]{
    \setlength{\anchoFloat}{#1\textwidth}
    \addtolength{\anchoFloat}{-4\fboxsep}
    \setlength{\anchoFigura}{\anchoFloat}
    \begin{figure}[#6]
        \begin{center}
            \Ovalbox{
                \begin{minipage}{\anchoFloat}
                    \begin{center}
                        \includegraphics[width=\anchoFigura,#5]{#2}
                        \caption{#3}
                        \label{#4}
                    \end{center}
                \end{minipage}
            }
        \end{center}
    \end{figure}
}

% Comando para incluir imágenes en formato apaisado (sin marco).
\newcommand{\figuraApaisadaSinMarco}[5]{%
    \begin{figure}%
        \begin{center}%
            \includegraphics[angle=90,height=#1\textheight,#5]{#2}%
            \caption{#3}%
            \label{#4}%
        \end{center}%
    \end{figure}%
}

% Para las tablas
\newcommand{\otoprule}{\midrule [\heavyrulewidth]}

% Nuevo comando para tablas pequeñas (menos de una página).
\newcommand{\tablaSmall}[5]{%
    \begin{table}
        \begin{center}
            \rowcolors {2}{gray!35}{}
            \begin{tabular}{#2}
                \toprule
                #4
                \otoprule
                #5
                \bottomrule
            \end{tabular}
            \caption{#1}
            \label{tabla:#3}
        \end{center}
    \end{table}
}

% Nuevo comando para tablas pequeñas (menos de una página).
\newcommand{\tablaSmallSinColores}[5]{%
    \begin{table}[H]
        \begin{center}
            \begin{tabular}{#2}
                \toprule
                #4
                \otoprule
                #5
                \bottomrule
            \end{tabular}
            \caption{#1}
            \label{tabla:#3}
        \end{center}
    \end{table}
}

\newcommand{\tablaApaisadaSmall}[5]{%
    \begin{landscape}
        \begin{table}
            \begin{center}
                \rowcolors {2}{gray!35}{}
                \begin{tabular}{#2}
                    \toprule
                    #4
                    \otoprule
                    #5
                    \bottomrule
                \end{tabular}
                \caption{#1}
                \label{tabla:#3}
            \end{center}
        \end{table}
    \end{landscape}
}

% Nuevo comando para tablas grandes con cabecera y filas alternas coloreadas en
% gris.
\newcommand{\tabla}[6]{%
    \begin{center}
        \tablefirsthead{
            \toprule
            #5
            \otoprule
        }
        \tablehead{
            \multicolumn{#3}{l}{\small\sl continúa desde la página anterior}\\
            \toprule
            #5
            \otoprule
        }
        \tabletail{
            \hline
            \multicolumn{#3}{r}{\small\sl continúa en la página siguiente}\\
        }
        \tablelasttail{
            \hline
        }
        \bottomcaption{#1}
        \rowcolors {2}{gray!35}{}
        \begin{xtabular}{#2}
            #6
            \bottomrule
        \end{xtabular}
        \label{tabla:#4}
    \end{center}
}

% Nuevo comando para tablas grandes con cabecera.
\newcommand{\tablaSinColores}[6]{%
    \begin{center}
        \tablefirsthead{
            \toprule
            #5
            \otoprule
        }
        \tablehead{
            \multicolumn{#3}{l}{\small\sl continúa desde la página anterior}\\
            \toprule
            #5
            \otoprule
        }
        \tabletail{
            \hline
            \multicolumn{#3}{r}{\small\sl continúa en la página siguiente}\\
        }
        \tablelasttail{
            \hline
        }
        \bottomcaption{#1}
        \begin{xtabular}{#2}
            #6
            \bottomrule
        \end{xtabular}
        \label{tabla:#4}
    \end{center}
}

% Nuevo comando para tablas grandes sin cabecera.
\newcommand{\tablaSinCabecera}[5]{%
    \begin{center}
        \tablefirsthead{
            \toprule
        }
        \tablehead{
            \multicolumn{#3}{l}{\small\sl continúa desde la página anterior}\\
            \hline
        }
        \tabletail{
            \hline
            \multicolumn{#3}{r}{\small\sl continúa en la página siguiente}\\
        }
        \tablelasttail{
            \hline
        }
        \bottomcaption{#1}
        \begin{xtabular}{#2}
            #5
            \bottomrule
        \end{xtabular}
        \label{tabla:#4}
    \end{center}
}

\definecolor{cgoLight}{HTML}{EEEEEE}
\definecolor{cgoExtralight}{HTML}{FFFFFF}

% Nuevo comando para tablas grandes sin cabecera.
\newcommand{\tablaSinCabeceraConBandas}[5]{%
    \begin{center}
        \tablefirsthead{
            \toprule
        }
        \tablehead{
            \multicolumn{#3}{l}{\small\sl continúa desde la página anterior}\\
            \hline
        }
        \tabletail{
            \hline
            \multicolumn{#3}{r}{\small\sl continúa en la página siguiente}\\
        }
        \tablelasttail{
            \hline
        }
        \bottomcaption{#1}
        \rowcolors[]{1}{cgoExtralight}{cgoLight}

        \begin{xtabular}{#2}
            #5
            \bottomrule
        \end{xtabular}
        \label{tabla:#4}
    \end{center}
}

% Capítulo
\newcommand{\capitulo}[1]{
    \stepcounter{chapter}
    \setcounter{section}{0}
    \setcounter{figure}{0}
    \setcounter{table}{0}
    \chapter*{\thechapter.\enskip #1}
    \addcontentsline{toc}{chapter}{\thechapter.\enskip #1}
    \markboth{#1}{#1}
}

\begin{document}

\renewcommand\maketitle{
    \thispagestyle{empty}

    \cabecera
    \vfill

    \begin{minipage}{0.74\textwidth}
        \begin{tcolorbox}[
                colback=blue!10,
                colframe=blue!20!black!40,
                top=16pt,
                bottom=16pt
            ]
            \begin{center}
                \large\textbf{TFG del Grado en Ingeniería Informática}

                \bigbreak

                \LARGE\textbf{\thetitle}
            \end{center}
        \end{tcolorbox}
    \end{minipage}
    \hfill
    \begin{minipage}{0.2\textwidth}
        \includegraphics[width=\textwidth]{escudoInfor}
    \end{minipage}


    \vfill

    \begin{center}
        \noindent\LARGE
        Presentado por \theauthor\\
        en Universidad de Burgos --- \thedate\\
        Tutor: \thetutor\\
    \end{center}

    \clearpage
}
\maketitle


\blankpage
\blankpage

% Tutor
\thispagestyle{empty}
\cabecera

\vfill

\nonzeroparskip{

    \noindent Dña. Alicia Olivares Gil y D. Álvar Arnaiz González, profesores
    del departamento de Ingeniería Informática, área de Lenguajes y Sistemas
    Informáticos.

    \noindent Expone:

    \noindent Que el alumno D. \theauthor, con DNI \thedni, ha realizado el
    Trabajo final de Grado en Ingeniería Informática titulado Identificación de
    Parkinson por visión artificial.

    \noindent Y que dicho trabajo ha sido realizado por el alumno bajo la
    dirección del que suscribe, en virtud de lo cual se autoriza su presentación
    y defensa.

}
\begin{center}
    En Burgos, a {\large \today}
\end{center}

\vfill\vfill\vfill

% Author and supervisor
\begin{minipage}{0.45\textwidth}
    \begin{flushleft}
        Vº. Bº. de la tutora:\\[2cm]
        Dña. Alicia Olivares Gil
    \end{flushleft}
\end{minipage}
\hfill
\begin{minipage}{0.45\textwidth}
    \begin{flushleft}
        Vº. Bº. del cotutor:\\[2cm]
        D. Álvar Arnaiz González
    \end{flushleft}
\end{minipage}

\hfill
\vfill


\blankpage
\blankpage


\frontmatter


% Resumen
\renewcommand*\abstractname{Resumen}
\begin{abstract}
    En este primer apartado se hace una \textbf{breve} presentación del tema que
    se aborda en el proyecto.
\end{abstract}
\renewcommand*\abstractname{Descriptores}
\begin{abstract}
    Palabras separadas por comas que identifiquen el contenido del proyecto Ej:
    servidor web, buscador de vuelos, android \ldots
\end{abstract}

\clearpage

% Resumen en inglés
\renewcommand*\abstractname{Abstract}
\begin{abstract}
    A \textbf{brief} presentation of the topic addressed in the project.
\end{abstract}
\renewcommand*\abstractname{Keywords}
\begin{abstract}
    keywords separated by commas.
\end{abstract}


\clearpage
\tableofcontents

\clearpage
\listoffigures

\clearpage
\listoftables



\mainmatter


\capitulo{Introducción}
\label{cha:Introducción}

La enfermedad de Parkinson es una enfermedad neurodegenerativa que afecta a
millones de personas en todo el mundo. Se caracteriza por la degradación de las
células nerviosas que controlan el movimiento, lo que puede provocar temblores,
rigidez muscular y pérdida del control postural \cite{eswiki:148845196}. Aunque
no existe una cura para la enfermedad de Parkinson, la detección temprana y el
tratamiento adecuado pueden mejorar significativamente la calidad de vida de los
pacientes.

En los últimos años, la inteligencia artificial se ha convertido en una
herramienta de gran utilidad para el diagnóstico y tratamiento de enfermedades,
incluida la enfermedad de Parkinson. En este proyecto, se ha desarrollado un
sistema basado en inteligencia artificial para crear modelos de aprendizaje
automático capaces de detectar si un individuo padece la enfermedad de Parkinson
a partir de un vídeo de la persona realizando un gesto característico con la
mano, además de cierta información adicional sobre el individuo.

Interactuar con los modelos creados de forma directa puede resultar complicado
si no se dispone de ciertos conocimientos previos sobre programación e
inteligencia artificial. Con el objetivo de facilitar esta tarea para la persona
media se ha desarrollado una aplicación web que simplifica el uso de los
modelos.

Los resultados obtenidos indican que los modelos creados son capaces de detectar
la enfermedad de Parkinson con bastante precisión, lo que sugiere que podrían
ser de utilidad para la detección de la enfermedad. Además, el enfoque basado en
vídeo es no invasivo y fácil de usar, lo que lo hace potencialmente útil en
entornos clínicos y de telemedicina.

Por último, me gustaría dar las gracias a todos aquellos que han dedicado su
tiempo y esfuerzo a la obtención de los vídeos utilizados para entrenar los
modelos, incluyendo la Asociación de Parkinson de Burgos, el personal del
Hospital Universitario de Burgos y el equipo de la Dra. Esther Cubo.

\capitulo{Objetivos del proyecto}
\label{cha:Objetivos del proyecto}

\section{Objetivos generales}

El objetivo principal de este proyecto es la creación de un sistema basado en la
inteligencia artificial que permita la detección de bradicinesia de forma
consistente utilizando como entrada un vídeo de la prueba de <<finger-tapping>>
(además de otra información), además este sistema debería ser utilizable por la
persona media. Para alcanzar este objetivo se plantean los siguientes objetivos
específicos:

\begin{enumerate}
    \item Revisar los trabajos previos relacionados para obtener unos
    conocimientos iniciales sobre el tema e identificar qué métodos funcionan
    mejor y peor.
    \item Diseñar e implementar una secuencia de pasos o <<pipeline>> para
    extraer características relevantes de los vídeos para su uso en algoritmos
    de aprendizaje automático.
    \item Buscar un modelo que se ajuste considerablemente bien a los datos
    disponibles mediante alguna de las técnicas existentes en el campo del
    aprendizaje automático para este propósito.
    \item Crear una aplicación web que permita subir un vídeo de la prueba,
    extraer una serie de características de este y proporcionárselas a un modelo
    de aprendizaje automático que sea capaz de devolver una predicción relativa
    a la presencia o no de bradicinesia en el paciente.
    \item Debido al público de una aplicación de este tipo se debe tener muy en
    cuenta la accesibilidad a la hora de diseñarla.
    \item Incluir en la aplicación web la funcionalidad de gestión de usuarios,
    así como la posibilidad de actualizar el modelo utilizado para realizar las
    predicciones por parte de un administrador.
    \item Documentar el trabajo realizado.
\end{enumerate}

\section{Objetivos técnicos}

Los siguientes objetivos han están relacionados con los aspectos técnicos de la
aplicación.

\begin{itemize}
    \item Utilizar Docker para facilitar el despliegue de la
    aplicación en cualquier equipo que tenga el motor de Docker instalado.
    \item Crear una API (Application Programming Interface) accesible mediante
    internet desde cualquier máquina para separar de forma clara el <<backend>>
    y el <<frontend>> y permitir la extensión futura de la aplicación.
    \item Internacionalizar la aplicación para permitir el uso de esta en varios
    idiomas y realizar una detección automática del idioma preferido por el
    usuario.
\end{itemize}
\capitulo{Conceptos teóricos}
\label{cha:Conceptos teóricos}

\capitulo{Técnicas y herraminetas}
\label{cha:Técnicas y herraminetas}

\section{Herramientas de desarrollo}

\subsection{Windows Subsystem for Linux}

Windows Subsystem for Linux (WSL) es una utilidad creada principalmente por
Microsoft y Canonical(Mantenedores de Ubuntu) que permite utilizar un sistema
operativo Linux desde Windows sin grandes pérdidas de rendimiento. Existen
varias opciones de distribución que se pueden utilizar (Kali Linux, OpenSUSE,
etc.), pero la mas utilizada y mejor mantenida para WSL es Ubuntu, por lo que es
la que se ha utilizado.

La mayor parte del desarrollo ha sido realizada desde WSL, donde se ha instalado
tanto \LaTeX como Python 3.9 y algunas dependencias, proceso mucho más simple de
realizar en Linux que en Windows.

\subsection{Visual Studio Code}

Visual Studio Code (VS Code) es un popular editor de texto altamente extensible
orientado al desarrollo creado por Microsoft. Gracias a su gran extensibilidad
es posible utilizarlo cómodamente en casi cualquier escenario. Mediante una de
estas extensiones (denominada literalmente WSL) se puede conectar con WSL y
utilizar los programas que tiene instalados como si hubiese sido ejecutado
diréctamente desde un sistema operativo Linux.

En este caso se ha utilizado VS Code para:

\begin{itemize}
    \item Crear tanto esta memoria como los anexos mediante la extensión
          \href{https://github.com/James-Yu/LaTeX-Workshop}{\LaTeX Workshop} que
          integra distintos compiladores de \LaTeX (pdfTeX, XeTeX, LuaTeX, etc.)
          dentro del programa y añade otras funcionalidades que hacen la edición
          de código \LaTeX mucho más conveniente.
    \item Crear la aplicación web, que utiliza el framework de JavaScript
          SvelteKit, mediante multitud de extensiones para integrar la multitud
          de librerías utilizadas (Svelte, TypeScript, Prettier, etc.).
    \item Editar archivos de configuración utilizados para la creación de los
          contenedores de Docker que componen la totalidad de la aplicación web
          (web, API y proxy inverso).
\end{itemize}

\subsection{PyCharm}

PyCharm es un Entorno de Desarrollo Integrado (IDE, por sus siglas en inglés)
creado y mantenido por JetBrains, una compañía que se dedica a crear software y
cuyos productos más representativos son multitud de IDEs para diferentes
lenguajes de programación (IntelliJ(Java), CLion(C y C++), GoLand(Go), etc.).

El objetivo de PyCharm es facilitar el desarrollo de aplicaciones que utilizad
Python y da soporte amplio tanto para los habituales ficheros \texttt{.py} como
para \texttt{.ipynb} (Notebooks de Jupyter), que han sido utilizados
extensamente en el proyecto.

Además, PyCharm por defecto permite la integración con entornos WSL de forma muy
fácil, simplemente se le debe indicar la localización del ejecutable de Python
en el sistema de archivos de la distribución de WSL.

En este proyecto PyCharm se ha utilizado para:

\begin{enumerate}
    \item Desarrollar la librería PADDEL, que contiene el código de Python común
    a los Notebooks de Jupyter utilizados durante la fase de investigación y a
    la aplicación web.
    \item Desarrollar la API que utiliza la aplicación web.
\end{enumerate}

\subsection{Git}

Actualmente Git es el sistema de control de versiones más extendido, en este
proyecto se ha usado para gestionar y guardar de forma segura los cambios que se
han ido realizado con el tiempo en la plataforma GitHub.

Aunque existen varias herramientas con interfaces gráficas que envuelven el
funcionamiento de Git (GitKraken, GitHub Desktop, etc.) para este proyecto se ha
utilizado diréctamente el comando \texttt{git} desde una terminal de WSL.

\capitulo{Aspectos relevantes del desarrollo del proyecto}
\label{cha:Aspectos relevantes del desarrollo del proyecto}

Este apartado contiene algunos aspectos relevantes sobre el desarrollo del
proyecto, incluyendo las razones detrás de las decisiones tomadas durante este y
el impacto que dichas decisiones han tenido sobre los resultados obtenidos.

\section{Fase de experimentación}

El conjunto de datos disponible durante la fase de experimentación consisten en
158 muestras con los siguientes ``atributos'':

\begin{itemize}
    \item Si la persona correspondiente a la instancia padece la enfermedad de
    Parkinson. Es lo que se busca predecir.
    \item Archivo de vídeo de una mano de la persona realizando la prueba de
    <<finger-tapping>>.
    \item Mano que aparece en el vídeo (izquierda o derecha).
    \item Fecha en la que se ha tomado el vídeo (irrelevante).
    \item Edad de la persona.
    \item Sexo de la persona.
    \item Mano dominante de la persona.
\end{itemize}

Cabe destacar que se ha tomado una muestra de cada mano de cada persona. Por lo
que este conjunto de datos ha sido obtenido de 79 personas. Aún así, se ha
tomado cada mano como una instancia distinta con el objetivo de tener más
instancias con las que trabajar.

Además, se dispone de 69 instancias de personas que padecen la enfermedad de
Parkinson y 89 instancias de personas sanas, por lo que se está trabajando con
un conjunto de datos ligeramente desbalanceado. Esto deberá tenerse en cuenta
para que no influya en los resultados obtenidos, por ejemplo, seleccionando
métricas adecuadas.

\subsection{Procesado de vídeo}

El objetivo de la fase de experimentación es crear un sistema que tenga por
entrada los atributos antes mencionados (salvo el objetivo) y por salida si esos
atributos corresponden con una persona que padece la enfermedad de Parkinson
(posiblemente junto con el grado de confianza de la predicción).

Los algoritmos de aprendizaje automático típicamente están diseñados para
trabajar con atributos numéricos, por lo que los atributos disponibles deben ser
transformados durante el preprocesado. Esto se puede realizar en la mayoría de
casos mediante simples métodos de minería de datos, como, por ejemplo, la
codificación de variables categóricas. Pero hay un problema, el archivo de
vídeo.

En este caso, una imagen es una matriz bidimensional donde cada celda contiene
una 3-tupla representando los valores rojo, verde y azul del pixel con el que se
corresponde. Un vídeo es una secuencia de imágenes junto con información
adicional, es este caso sólamente es relevante la tasa de refresco, que es la
frecuencia a la que ha sido capturada cada imagen. Esto es de gran importancia
para dar una componente temporal al vídeo.

Con lo anterior se puede ver que un vídeo no se puede utilizar de forma directa,
sino que deberá pasar por una fase multietapa de preprocesado para ser
transformado a un conjunto de valores numéricos que lo describen.

\subsubsection{Extracción de <<landmarks>>}

Una <<landmark>> en el campo de la visión aritificial es un punto de referencia
que se corresponde con un objeto físico. El objetivo de esta fase es convertir
el vídeo de la prueba de <<finger-tapping>> a una secuencia o serie temporal de
conjuntos de <<landmarks>> los cuales representan la posición de la mano durante
un instante concreto.

El primer paso para lograr esto es leer las imágenes del archivo de vídeo, esto
ha sido logrado mediante la librería de <<bindings>> de OpenCV para Python.
OpenCV es una librería de uso general para realizar tareas relacionadas con la
visión artificial, y permite leer multitud de formatos de vídeo mediante un
flujo de imágenes.

Se ha tomado especial precaución para no cargar en memoria todas las imágenes
del vídeo, sino ir leyéndolas cuando son necesarias. Un archivo de vídeo ocupa
relativamente poco espacio en disco debido a que está codificado mediante un
<<codec>> determinado, lo que reduce el tamaño del archivo. Pero, para poder
trabajar con las imágenes del vídeo, este debe ser decodificado.

La magnitud de lo anterior se ve claramente con el siguiente ejemplo. Un vídeo
con una resolución de 1920x1080 a 30 fotogramas por segundo codificado mediante
el codec H264 que ocupa 50,7MB, decodificado, ocuparía en torno a los 3,7GB.

A continuación se deben extraer las <<landmarks>> de las imágenes leidas. Esto
ha sido logrado mediante la librería Mediapipe Hands \cite{zhang2020mediapipe},
que permite extraer los siguentes puntos a partir de una imágen:

\imagen{relevant_aspects/hand-landmarks.png}{Mediapipe Hands <<Landmarks>>}{1}

Debido a la perspectiva desde la que se han tomados los vídeos de la prueba de
<<finger-tapping>> hay algunas <<landmarks>> que quedan ocultas en varias
ocasiones, lo que hace su posición detectada poco fiable. Por lo que se desechan
las <<landmarks>> en el intervalo [9-17] para solo tomar en cuenta la parte
``frontal'' de la mano.

Mediapipe Hands ofrece la posibilidad de extraer <<landmarks>> de las imágenes
como si estas no tuviesen relación entre sí o realizar la extracción teniendo en
cuenta que se trata de un vídeo, así teniendo en cuenta resultados anteriores
para determinar de forma más precisa las <<landmarks>> del fotograma actual.
Esto tiene la gran desventaja de que hace imposible la paralelización de este
proceso (debido a que cada extracción depende de los resultados anteriores).

Lo anterior no tiene un gran efecto durante la fase de investigación, ya que sí
se puede paralelizar la extracción de <<landmarks>> de cada vídeo. Donde sí
tiene efecto es sobre la aplicación web, empeorando la experiencia de los
usuarios.

En última instancia se tomó la decisión de no usar paralelización, y así obtener
resultados más precisos. Una alternativa planteada fue realizar la extracción en
paralelo solamente en la aplicación, pero esto se descartó ya que esto podría
introducir inconsistencias y variables adicionales que deberían ser tomadas en
cuenta.

\subsubsection{Extracción de series temporales}

Tras el proceso descrito en el apartado anterior se dispone una secuencia de
datos similar a la siguiente:

\imagen{relevant_aspects/poses_series.png}{Serie temporal de <<landmarks>>}{1}

Cada fila representa la pose de la mano durante un fotograma concreto. Se puede
observar que el índice es una fecha completa en lugar de otras opciones más
convenientes como los segundos desde el inicio del vídeo o lo que se conoce como
<<timestamp>>. Esto es una imposición de TSFresh para poder realizar la
extracción de características teniendo en cuenta esta componente temporal.

Cada celda representa el punto en el espacio en el que se encuentra la
<<landmark>> correspondiente en un instante de tiempo.

Aunque estas series temporales son una simplificación muy grande de lo vídeos
iniciales, aún no son del todo útiles. Al ser información sobre la posición en
el espacio, es dependiente de la forma en la que ha sido tomada la grabación. En
un entorno controlado esto no sería mucho problema, pero en este caso, como se
desea hacer el uso del sistema accesible a cualquier persona sin necesitar
equipamiento especial, se deben buscar alternativas.

En este caso se ha decidido tomar el ángulo entre el ángulo entre la punta del
dedo pulgar, la muñeca y la punta del dedo índice. Esta medida es útil ya que no
depende de cómo se ha tomado el vídeo ni del tamaño de la mano de la persona.
Quedando como resultado una serie temporal como la siguiente:

\imagen{relevant_aspects/angle_series.png}{Serie temporal de ángulos}{0.5}

La unidad usada en estos ángulos es el radián.

\subsection{Extracción de características}



\capitulo{Trabajos relacionados}\label{cha:Trabajos relacionados}

Durante la última década se ha intentado utilizar la visión por computador para
evaluar la Enfermedad de Parkinson en múltiples ocasiones con diferentes
resultados. Este capítulo recopila de forma resumida algunos de los trabajos más
relevantes.


\section{A computer vision framework for finger-tapping evaluation}

Este artículo~\cite{khan2014computer} documenta el uso de visión por computador
para determinar el nivel de severidad de la Enfermedad de Parkinson y distinguir
entre individuos con esta enfermedad e individuos sin ella.

El método empleado se caracteriza por utilizar vídeos con la cara de la persona
y ambas manos a los lados de la cabeza, apuntando las puntas de los dedos hacia
la misma. Esto se utiliza para poder normalizar las distancias en base a
características faciales.

El estudio se realizó sobre 13 pacientes con la Enfermedad de Parkinson, tomando
17 vídeos de cada paciente durante un día, además de un grupo de control de 6
individuos, tomando 2 vídeos al día durante una semana por cada uno. Aunque
algunos de estos vídeos fueron descartados. En total se utilizaron 471 vídeos.


\paragraph{Metodología}

\begin{enumerate}
    \item Se detecta la cara del individuo para la normalización. Esto se basa
          en que la longitud de la mano de una persona adulta es aproximadamente
          igual a la altura de su cara.
    \item Se obtiene una serie temporal que representa la amplitud del
          movimiento de los dedos índice y pulgar de la mano dominante del
          individuo.
    \item Se extraen un total de 15 características de esta serie temporal, por
          ejemplo:
          \begin{itemize}
              \item Correlación cruzada media entre los máximos locales de dos
                    intervalos distintos de tiempo de la serie temporal. Esto
                    mismo se realiza también sobre los mínimos locales.
              \item Número total de toques de dedos durante la grabación.
              \item Velocidad media de la apertura de dedos.
              \item Velocidad media del cierre de dedos.
              \item \dots
          \end{itemize}
    \item Se realiza una selección de características eliminando aquellas
          redundantes y usando el algoritmo chi-cuadrado.
    \item Se entrena una máquina de vectores de soporte (SVM) mediante las
          características obtenidas para realizar la clasificación.
\end{enumerate}


\paragraph{Resultados}

En cuanto a la distinción entre pacientes de la Enfermedad de Parkinson y el
grupo de control se obtuvo una precisión del 95\%, que es una cifra que se
debería tomar con precaución debido que, aunque se han utilizado 471 vídeos,
estos provienen de únicamente 19 personas.


\section{The discerning eye of computer vision}

En este estudio \cite{williams2020discerning} realizado sobre 39 pacientes con
la Enfermedad de Parkinson y sobre un grupo de control de 30 individuos se
tomaron vídeos de ambas manos de cada individuo mientras realizan toques de los
dedos índice y pulgar (de forma similar a cómo se han tomado las muestras para
este trabajo). Dando un total de 133 vídeos (se descartó uno).

De estos vídeos se han extraído diferentes características y comprobado la
relación que existe entre éstas y diferentes escalas que clasifican el nivel de
gravedad de la Enfermedad en un paciente, como la \textit{Modified Bradykinesia
    Rating Scale} (MBRS) que categoriza el movimiento de los individuos en 5
niveles, del 0 al 4, siendo 0 un movimiento normal y 4 el nivel de mayor
gravedad.


\paragraph{Metodología}

\begin{enumerate}
    \item Se utiliza una librería de visión por computador, en concreto
          DeepCutLab, para obtener una serie temporal de la amplitud entre las
          puntas de los dedos pulgar e índice.
    \item Se normaliza esta serie temporal utilizando la amplitud máxima
          detectada, que va a convertirse en el valor 1, siendo todos los demás
          valores escalados proporcionalmente.
    \item Se extraen las siguientes características:
          \begin{itemize}
              \item Velocidad, calculada como la tasa media de cambio.
              \item Variabilidad de la amplitud, calculada como el coeficiente
                    de variación de la diferencia media entre máximos y mínimos
                    de diferentes intervalos de 1 segundo de la serie temporal.
              \item Regularidad del ritmo, calculada utilizando la Transformada
                    Rápida de Fourier y, a continuación, midiendo la potencia de
                    la frecuencia dominante más la potencia de las frecuencias
                    en un intervalo de 0.4 Hz alrededor de ésta (un ritmo más
                    regular concentra una mayor potencia en una única
                    frecuencia).
          \end{itemize}
\end{enumerate}


\paragraph{Resultados}

Se observó una correlación bastante alta entre las características utilizadas y
la categoría del individuo dentro de las escalas de medición de la Enfermedad de
Parkinson utilizadas medida por un experto en el campo.


\section{Supervised classification of bradykinesia}

Este estudio \cite{williams2020supervised} es muy similar al anteriormente
explicado, y está realizado por un equipo compuesto por casi los mismos
participantes. En este caso se utilizaron 70 vídeos, de ambas manos de 20
pacientes con la Enfermedad de Parkinson y de un grupo de control de 15
individuos.


\paragraph{Metodología}

La metodología es prácticamente igual que antes, la diferencia principal está en
las características que se extraen de la serie temporal correspondiente con la
amplitud, se ha obtenido:

\begin{itemize}
    \item Frecuencia, medida como la frecuencia máxima de la Transformada Rápida
          de Fourier de la serie temporal.
    \item Amplitud, calculada como la densidad espectral, que se ha obtenido
          mediante la integral cuadrada del espectro de la Transformada Rápida
          de Fourier.
\end{itemize}

Con estas características se ha realizado clasificación binaria mediante
clasificación bayesiana ingenua (naive bayes), regresión logística y máquina de
vectores de soporte, tanto con función lineal como con función de base radial.


\paragraph{Resultados}

Los mejores resultados se obtuvieron con máquina de vectores de soporte con
función de base radial, que coincide en un 73\% de los casos con la
clasificación de expertos en el campo.

\capitulo{Conclusiones y Líneas de trabajo futuras}
\label{cha:Conclusiones y Líneas de trabajo futuras}

Por último esta sección recopila las conclusiones a las que se ha llegado
durante la realización del proyecto. Además de algunas líneas de trabajo
futuras.

\section{Conclusiones}

Se han extraido las siguientes conclusiones durante el proyecto.

\begin{itemize}
    \item En el campo de la minería de datos la fase de preprocesado tiene tanta
    importancia como el propio entrenamiento de los modelos. En esta fase se
    determina lo complejos que deberán ser los modelos para adaptarse a los
    datos y el tiempo de entrenamiento que será necesario.
    \item Se ha creado una librería de Python instalable mediante \texttt{pip}
    para centralizar el código común la fase de investigación (generación de
    modelos) y la de desarrollo (aplicación web). Esta librería contiene todo el
    código relacionado con el procesado de los vídeos para transformarlos a un
    conjunto de características.
    \item Para facilitar la interacción con los modelos generados por cualquier
    programa, siempre que pueda realizar peticiones HTTP, se ha desarrollado una
    API Rest mediante el framework FastAPI.
    \item Se ha creado una aplicación web mediante un framework de JavaScript
    (SvelteKit) para permitir a cualquier persona que disponga de una navegador
    web interactuar con la API.
    \item Se ha utilizado Docker para crear contenedores con las distintas
    partes que componen la aplicación, simplificando en gran medida el proceso
    de configuración del entorno de desarrollo y producción.
    \item Con el objetivo de agilizar el proceso de creación de modelos se ha
    implementado un <<script>> de Python que presenta una interfaz por línea de
    comandos mediante la que se determinan distintos parámetros para crear un
    modelo que se puede subir a la aplicación web.
\end{itemize}

\section{Lineas de trabajo futuras}

Algunas líneas de trabajo futuro a corto o medio plazo:

\begin{itemize}
    \item Buscar alguna forma para acelerar el procesado de los vídeos,
    actualmente esto toma unos 20 segundos, que es una espera muy larga al
    utilizar la aplicación web. Una idea que no se llegó a implementar es el uso
    de programación concurrente para <<aprovechar>> el tiempo de espera de
    entrada/salida existente cuando se leen los fotogramas del disco.
    \item Añadir la posibilidad de entrenar modelos directamente desde la
    aplicación web, eliminando la necesidad de configurar el entorno para esto
    de forma local.
    \item Permitir el uso de múltiples modelos distintos desde el formulario
    para obtener predicciones, posiblemente como un menú de <<Opciones
    avanzadas>>.
    \item Añadir algún método de actualización de token de sesión JWT para
    impedir que el usuario pierda su sesión mientras usa la aplicación web.
\end{itemize}

El conjunto de datos utilizado contine muestras de personas que ya se conoce que
padecen la enfermedad de Parkinson, debido a esto, en este proyecto se ha
conseguido detectar la presencia de la enfermedad una vez está en un estado
avanzado. Esto no es muy útil para realizar una detección temprana y empezar un
tratamiento más efectivo en mermar los efectos de la enfermedad.

Una línea de trabajo futuro a largo plazo podría consistir en tomar vídeos de la
prueba de <<finger-tapping>> de forma aleatoria sobre la población, si la
muestra es lo suficientemente grande, algunas de esas personas desarrollarán la
enfermedad. Con esto podría ser posible crear modelos capaces de detectar
personas que van a padecer la enfermedad de Parkinson en el futuro, pudiendo
tomar las medidas precautorias necesarias con antelación.



\bibliographystyle{plain}
\bibliography{refs}

\end{document}
